\documentclass[12pt]{article}
\usepackage[utf8]{inputenc}
\usepackage[T1]{fontenc}
\usepackage{hyperref}
\title{Merwari Language Resources for K-12 Curriculum Development}
\author{Compiled by Claude Sonnet 5.5 (claude-sonnet-5-5) for Andr\'as Kornai}
\date{}
\begin{document}
\maketitle

\section{Introduction}
Merwari (ISO 639-3 wry; Glottocode merw1238) is a Western Indo-Aryan variety of India within the Marwari macrolanguage and the Rajasthani group. It is spoken in the Ajmer and Nagaur districts of Rajasthan (as described in the Wikipedia summary), and its educated speakers are also fluent in Hindi. The searches found no resource specific to Merwari. Most useful material is for Marwari generally.

\section{Primary Educational Resources}
No Merwari primers, textbooks or school programmes were found.

\section{Linguistic and Reference Materials}
\begin{itemize}
\item Wikipedia, Marwari language, which covers Merwari as a variety: \url{https://en.wikipedia.org/wiki/Marwari_language}
\item Webonary Marwari glossary. It is a work in progress and describes itself as far from complete. It is Marwari in general, not Merwari specifically: \url{https://www.webonary.work/marwari/?p=9444}
\end{itemize}
No Merwari-specific grammar or dictionary was found.

\section{Religious and Community Resources}
None were identified in the searches.

\section{Digital and Audio Resources}
None specific to Merwari were found.

\section{Curriculum Development Context}
Merwari has a large speaker base, but I found no evidence of use in schools or of dedicated materials. Any curriculum would likely have to draw on Marwari and Rajasthani resources and on Hindi, the language of schooling. Those resources were not examined here.
\end{document}